% =====================================================================
%  1.6 Modelo de gestão  (Atividade 3.2.2.1)
%  1.6.1 Trabalho técnico-social  (Atividade 3.2.2.2)
% =====================================================================
\section{Modelo de gestão}
\label{subsec:modelos-gestao}

O modelo de gestão garante que os sistemas continuem funcionando depois da
entrega. Estado, municípios e comunidades compartilham responsabilidades. O
IAT, a SECID e o IDR-Paraná apoiam o planejamento e o monitoramento, os
municípios participam da governança, e as comunidades, por meio de associações
regularizadas, operam os sistemas. O modelo está em definição, e a alternativa
em análise é um arranjo intermunicipal, como um consórcio público
\cite{mop2026}. A adesão começa pelo município e chega à comunidade
(\autoref{fig:adesao}). A implantação segue as etapas da
\autoref{fig:modelo-gestao}, detalhadas no \autoref{qua:etapas-3221}
(Apêndice~\ref{ap:quadros}).

\begin{figure}[htbp]
  \centering
  \caption{Seleção e adesão dos municípios e comunidades, por responsável}
  \label{fig:adesao}
  \fluxograma{fluxograma-adesao}
  \fonte{Elaborado pelo IAT (2026) com base em \citeonline{mop2026} e no POA 2026--2027.}
\end{figure}

\begin{figure}[htbp]
  \centering
  \caption{Etapas de implantação do modelo de gestão e seus produtos}
  \label{fig:modelo-gestao}
  \fluxograma{fluxograma-modelo-gestao}
  \fonte{Elaborado pelo IAT (2026) com base em \citeonline[Quadro 28]{mop2026}.}
\end{figure}

\pendencia{A planilha de custos prevê três ``unidades'' de implantação do
modelo de gestão, de R\$ 2,3 milhões cada, sem definir o que é uma unidade.
No Quadro 28 do MOP, a etapa ``Supervisão técnica'' repete as evidências das
obras, e a sigla ``EC'' não está definida. Definir e harmonizar.}

% ---------------------------------------------------------------------
\subsection{Trabalho técnico-social}
\label{subsec:tts}

O trabalho técnico-social é executado por empresa contratada. Mobiliza a
comunidade, cadastra as famílias por busca ativa, faz o diagnóstico de água e
esgoto de cada domicílio e colhe o termo de adesão das famílias enquadradas
(\autoref{fig:tts}). O diagnóstico orienta o projeto do sistema de água e
define os domicílios da frente de esgotamento nas comunidades. A execução em
campo começa em outubro de 2027 (POA).

\begin{figure}[htbp]
  \centering
  \caption{Sequência e produtos do trabalho técnico-social}
  \label{fig:tts}
  \fluxograma{fluxograma-tts}
  \fonte{Elaborado pelo IAT (2026) com base em \citeonline{mop2026} e no POA 2026--2027.}
\end{figure}

\pendencia{O POA reúne em uma só contratação o trabalho técnico-social e os
estudos iniciais do poço, que o MOP trata separadamente (Atividades 3.2.2.2 e
3.2.2.3). Harmonizar.}
